%% kapa: exam=1 round=24 year=2013 period=1 form=A subject=civil_law subject_ko=민법 seq=14 number=14
%% source: Q-Net 감정평가사 기출문제 (cst00309) artlSeq=5073802 file=extracted/1008/5073802/2013년 제24회 감정평가사 1차시험 기출문제/1교시(A형).hwp
%% conversion: kordoc 4.24.51 -> tools/gen1.py (auto); review=proofread-sonnet; flags=none
\begin{question}{14}
乙이 대리권 없이 甲의 대리인으로 甲 소유의 토지를 丙에게 임대한 경우에 관한 설명으로 옳지 않은 것은? (다툼이 있으면 판례에 의함)
\begin{choices}
  \item 甲이 계약을 추인한 경우 다른 의사표시가 없으면 그 계약은 처음부터 유효하나, 제3자의 권리를 해치지 못한다.
  \item 乙이 이행책임 또는 손해배상책임을 져야 하는 경우 어느 책임을 부담할 것인지는 원칙적으로 丙이 선택한다.
  \item 甲이 사망하고 乙이 유일한 상속인인 경우 乙은 이 계약의 추인을 거절하지 못한다.
  \item 丙이 상당한 기간을 정해 추인 여부를 최고한 경우 甲이 그 기간 내에 확답을 발송하지 않으면 추인을 거절한 것으로 본다.
  \item 丙이 선의인 경우 그는 甲이 계약을 추인한 후에도 그 계약을 철회할 수 있다.
\end{choices}
\end{question}
\answer{5}
